\section{Conclusion}  

Adaptation to new information across a range of timescales is a crucial ability for generally intelligent agents. Foundation models in particular have demonstrated an ability to acquire a large knowledge-base of information, and apply this rapidly to new scenarios. Thus far, they have relied mainly on supervised and self-supervised learning. As such, they require access to large datasets. An alternative to collecting datasets is to have an agent learn from its own experience via reinforcement learning, provided that sufficiently rich physical worlds or open-ended simulations are available. This raises the question: can large-scale, generally adaptive models be trained with RL?

In this paper, we demonstrate, for the first time to our knowledge, an agent trained with RL that is capable of rapid in-context adaptation across a vast, open-ended task space, at a timescale that is similar to that of human players. This \textit{Adaptive Agent} (AdA) explores held-out tasks in a structured way, refining its policy towards optimal behaviour given only a few interactions with the task. Further, AdA is amenable to contextual first-person prompting, strengthening its few-shot performance, analogous to prompting in large language models. AdA shows scaleable performance as a function of number of parameters, context length and richness of the training task distribution. 

Our training method is based on black-box meta-RL, previously thought of as hard to scale. We show that state-of-the-art automatic curriculum techniques can shape the data distribution to provide sufficient signal for learning to learn in an open-ended task space. Moreover, we demonstrate that attention-based architectures can take advantage of this signal much more effectively than purely recurrent networks, illustrating the importance of co-adapting data-distribution and agent architecture for facilitating rapid adaptation. Finally, distillation enables us to realise the potential of large-scale Transformer architectures. 

The future of AI research will inevitably involve training increasingly large models with increasingly general and adaptive capabilities. In this direction, we have provided a recipe for training a 500M parameter model, which we hope can pave the way for further advances at the intersection of RL and foundation models. AdA shows rapid and scalable adaptation of myriad kinds, from tool use to experimentation, from division of labour to navigation. Given scaling law trends, such models may in future become the default foundations for few-shot adaptation and fine-tuning on useful control problems in the real world.